\documentclass[11pt]{scrartcl}

\usepackage[stixtwo]{fontsetup}
\usepackage{amsmath}

\begin{document}
$x^2 - x - 2 = (x - 2)(x + 1)$ és $x^2 - 1 = (x - 1)(x + 1)$.
\begin{enumerate}
	\item $x > 2$ esetén $|x^2 - x - 2| = x^2 - x - 2$, így
	      \[
		      \lim_{x \to +\infty} f(x)
		      = \lim_{x \to +\infty}
		      \frac{1 - \frac{1}{x} - \frac{2}{x^2}}{1 - \frac{1}{x^2}}
		      = 1
	      \]

	\item Ha $-1 < x < 1$, akkor $x - 2 < 0$ és $x + 1 > 0$, így
	      $|x^2 - x - 2| = (2 - x)(x + 1)$, tehát
	      \[
		      \lim_{x \to -1 + 0} f(x)
		      = \lim_{x \to -1 + 0} \frac{(2 - x)(x + 1)}{(x - 1)(x + 1)}
		      = \lim_{x \to -1 + 0} \frac{2 - x}{x - 1}
		      = \frac{3}{-2}
		      = -\frac{3}{2}
	      \]

	\item \[
		      \lim_{x \to 1 - 0} f(x)
		      = \lim_{x \to 1 - 0}
		      \frac{\overbrace{|x^2 - x - 2|}^{\to 2}}{\underbrace{x^2 - 1}_{\to 0 - 0}}
		      = -\infty
	      \]
\end{enumerate}
\end{document}
